The rule's reporting boundary has inadequate null behavior, and its fixed-pool
coverage estimates are 84.5--91.0\% at that boundary in the paired diagnostic.
Near-perfect inclusion of the sampled reference at larger counts overstates
coverage of the generating-pool effect. The variance calculation identifies
why changing the target changes this assessment. Separately, first-crossing
candidates can fail at larger counts, and cohort differences in relative
expression variance help explain the different detection curves.

The evidence establishes a failure in this single-cell application. A
replacement rule needs a specified target, joint assessment of null rejection,
effect detection and abstention, and independent evaluation after selection.
The synthetic external-control example shows why an independently constructed
reference must also average over the intended target population. The broader
patient-population question requires new patients and cohorts.
